\section{Open Challenges and Future Research Directions}
\label{sec:future}

\subsection{From capable legal AI to dependable legal information retrieval}

The preceding sections show that contemporary legal AI has made progress in legal language
understanding, evidence retrieval, document drafting, and agentic processing. However, these
capabilities remain difficult to integrate into a system that consistently returns relevant, applicable,
and verifiable legal answers.

The central research challenge is therefore not simply to produce more fluent legal text. It is to
develop systems that can retrieve authoritative evidence, preserve its legal context, reason over
relationships between sources, and verify the claims they generate. The following research directions
emerge from the comparison of legal AI architectures and the challenges of jurisdiction-specific
retrieval.

\subsection{Research direction 1: Knowledge-graph-grounded legal retrieval}

Conventional retrieval methods primarily rank text passages according to lexical or semantic relevance.
While effective for many questions, they may struggle with queries requiring relationships across
multiple documents. For example, answering a complex legal question may require connecting a statutory
provision to a definition, identifying judgments that interpret it, and determining whether a later
authority modifies or distinguishes an earlier interpretation. A flat index may retrieve individual
passages without explicitly representing these relationships.

Knowledge graphs provide a complementary representation by encoding entities and relationships as
structured edges. In legal applications, graph entities might represent statutes, provisions, judgments,
courts, parties, legal concepts, and dates. Relationships could represent references, interpretation,
amendment, procedural connections, or other explicitly supported links.

Future research should investigate how to combine graph-based retrieval with lexical and dense retrieval
while preserving provenance for every relationship. The graph must distinguish explicitly stated facts
from inferred connections and should not encode a legal relationship merely because two documents discuss
similar subjects. Key research questions include:
\begin{itemize}[leftmargin=1.4em, itemsep=2pt, topsep=2pt]
\item Does graph-grounded retrieval improve recall for multi-hop legal questions compared with hybrid
text retrieval alone?
\item Can graph relationships improve the identification of relevant authorities and exceptions?
\item How should uncertain, conflicting, or incomplete graph relationships be represented?
\item What methods can update the graph when statutes are amended or new judgments become available?
\end{itemize}

The most useful evaluation would compare a text-only baseline, a hybrid retrieval system, a graph-only
method, and a combined graph--text architecture using the same test queries and evidence annotations.

\subsection{Research direction 2: Multilingual and low-resource legal retrieval}

Multilingual legal retrieval remains challenging when queries are expressed in languages or scripts that
are poorly represented in training data. English, Urdu, and Roman Urdu may express similar legal
intentions through different lexical forms, transliterations, or terminology.

Future research should investigate cross-lingual retrieval models, legal terminology alignment,
transliteration-aware indexing, and query expansion. Rather than translating every query into English and
discarding the original wording, a system could search across multiple query representations and fuse the
results. The principal challenge is to distinguish linguistic similarity from legal equivalence. A
translated phrase may be grammatically correct but fail to represent a jurisdiction-specific legal concept
accurately. Useful research questions include:
\begin{itemize}[leftmargin=1.4em, itemsep=2pt, topsep=2pt]
\item How does retrieval performance vary between English, Urdu, and Roman Urdu queries?
\item Does multilingual dense retrieval outperform translation-only and transliteration-only baselines?
\item Can legal entity linking reduce errors involving statutory names, section numbers, courts, and
legal terminology?
\item How can multilingual benchmarks evaluate both retrieval relevance and legal correctness?
\end{itemize}

A reproducible benchmark should include queries in each supported language, authoritative source passages,
equivalent or carefully aligned legal intents, and explicit judgments of whether retrieved evidence is
applicable. It should also test spelling variation, code-switching, and unseen transliteration forms.

\subsection{Research direction 3: Structure-aware legal indexing}

Legal document structure is frequently lost during conversion, chunking, and indexing. This creates a
risk that the retrieval system will return a fragment without the definitions, exceptions, headings, or
parent sections needed to interpret it correctly.

Structure-aware indexing offers a way to preserve these relationships. Future research should compare
fixed-size chunking, semantic chunking, hierarchical indexing, and parent-child retrieval under the same
evaluation conditions. A particularly important question is whether preserving structure improves
downstream answer quality beyond its effect on retrieval recall. For example, a structure-aware system
may retrieve a relevant subsection and its parent heading, allowing the generator to interpret the
provision more accurately even if the retrieved subsection itself is unchanged.

Research should also examine how to handle malformed PDFs, inconsistent numbering, footnotes, tables,
nested clauses, and cross-references. These are not merely preprocessing inconveniences: errors during
document parsing can propagate through retrieval and generation. Promising evaluation measures include:
\begin{itemize}[leftmargin=1.4em, itemsep=2pt, topsep=2pt]
\item structural extraction accuracy for headings, sections, and subsections;
\item retrieval recall for provisions whose meaning depends on parent context;
\item accuracy of cross-reference resolution;
\item citation localization accuracy at the section or paragraph level;
\item downstream legal correctness under controlled comparisons.
\end{itemize}

\subsection{Research direction 4: Claim-level evidence verification}

Citation-grounded generation is insufficient if the system only checks whether a response contains
references. Future legal AI architectures should explicitly associate each material claim with the passage
or passages intended to support it.

A claim-level verification pipeline could identify factual and legal assertions, resolve their citations,
compare each claim with its evidence, and flag unsupported or contradictory statements. Depending on the
application, verification might combine deterministic citation resolution, metadata checks, textual
entailment models, and qualified human review. The research challenge is that legal support is often more
complex than ordinary textual entailment. A passage may establish a general rule while a different passage
provides an exception. A judgment may discuss an argument without adopting it. A later authority may change
how an earlier decision should be understood.

Future work should therefore evaluate more than whether a cited passage is topically relevant. It should
assess whether the passage supports the claim in context and whether material qualifications or
counter-authorities have been omitted. Possible research questions include:
\begin{itemize}[leftmargin=1.4em, itemsep=2pt, topsep=2pt]
\item How accurately can automated methods distinguish supported, contradicted, and unsupported legal
claims?
\item Can verification identify citations to the wrong subsection or an irrelevant paragraph within an
otherwise valid authority?
\item Does claim-level verification reduce unsupported-claim rates without introducing excessive false
alarms?
\item When should the system abstain rather than generate a legal conclusion?
\end{itemize}

An effective benchmark should include claim-level annotations, source locations, support labels, and
cases involving exceptions or conflicting authorities.

\subsection{Research direction 5: Agentic legal AI with explicit quality gates}

Agentic AI can coordinate multiple operations, such as query decomposition, iterative retrieval, legal
analysis, and citation checking. However, orchestration introduces additional opportunities for error. An
agent may select an inappropriate tool, carry forward a mistaken assumption, or produce a confident
conclusion based on incomplete evidence.

Future research should focus on evidence-controlled orchestration rather than simply increasing the number
of agents. Each stage should have explicit input requirements, output contracts, and quality checks. For
example, an agent responsible for legal analysis should receive retrieved evidence with source identifiers
and relevant document context. A verification stage should test whether the proposed claims are supported
before the final response is released. If evidence is missing, the system should trigger additional
retrieval or qualify its answer.

Software engineering principles are particularly relevant here. Modular interfaces, replaceable
components, explicit error handling, versioned configurations, and regression tests can make agentic legal
systems easier to maintain and evaluate. Potential experiments should compare a single-agent baseline with
multi-stage and multi-agent alternatives under equivalent retrieval and generation settings. Measurements
should include legal correctness, citation support, failure recovery, latency, computational cost, and the
rate of unsupported claims. Without such controls, improvements cannot reliably be attributed to agentic
orchestration.

\subsection{Research direction 6: Temporal and authority-aware reasoning}

Legal information changes over time. Statutes may be amended, judgments may be distinguished or overruled,
and procedural requirements may change. A retrieval system that ignores these changes can produce an
answer that was once correct but is no longer applicable.

Future systems should represent source dates, effective periods, amendments, and relevant relationships
between authorities. Where available and properly established, this information can support temporal
filtering and authority-aware retrieval. An important distinction is between the date a document was
published and the period during which its legal rule applied. These dates may differ. A reliable system
should not assume that the newest document is necessarily the most authoritative or applicable.

Research should evaluate whether temporal and authority-aware retrieval reduces outdated-law errors, and
how systems should respond when the corpus does not contain enough information to establish current legal
status.

\subsection{Research direction 7: Reproducible evaluation and realistic legal benchmarks}

Legal AI evaluation remains fragmented across tasks, jurisdictions, corpora, and definitions of
correctness. A model may perform well on a curated benchmark but struggle with ambiguous queries,
incomplete evidence, or questions that require multiple sources.

Future benchmarks should combine retrieval and generation evaluation while preserving separate scores for
each stage. They should include difficult cases involving exceptions, cross-references, competing
authorities, temporal changes, multilingual queries, and insufficient evidence. A strong evaluation
protocol should also report dataset construction, source coverage, annotation instructions, reviewer
expertise, inter-annotator agreement, baseline configurations, and uncertainty estimates. Ablation studies
should isolate the contributions of hybrid retrieval, rank fusion, reranking, hierarchical indexing,
knowledge graphs, and verification.

Public datasets and reproducible evaluation code can help researchers distinguish genuine architectural
improvements from differences in test construction or experimental settings. Where legal or privacy
restrictions prevent full data release, researchers should document the limitations and provide
reproducible alternatives where possible.

\subsection{Consolidated research agenda}

Table~\ref{tab:agenda} draws the directions together with measurable research objectives.

\begin{table*}[tp]
\centering
\caption{Open challenges and measurable research objectives.}
\label{tab:agenda}
\footnotesize
\setlength{\tabcolsep}{5pt}
\renewcommand{\arraystretch}{1.3}
\begin{tabularx}{\textwidth}{>{\raggedright\arraybackslash}p{3.4cm} >{\raggedright\arraybackslash}X >{\raggedright\arraybackslash}X}
\toprule
\textbf{Research challenge} & \textbf{Proposed direction} & \textbf{Example evaluation} \\
\midrule
Missing cross-document relationships & Combine knowledge graphs with text retrieval. & Recall on multi-hop questions and graph-relation accuracy. \\
Low-resource multilingual queries & Multilingual retrieval with terminology alignment. & Language-specific Recall@$k$ and expert-assessed correctness. \\
Loss of legal hierarchy & Structure-aware indexing and parent-child retrieval. & Structural extraction accuracy and downstream answer quality. \\
Unsupported legal claims & Claim-level evidence verification. & Claim-support precision, recall, and unsupported-claim rate. \\
Unreliable agentic workflows & Modular orchestration with explicit quality gates. & Correctness, failure recovery, latency, and ablation results. \\
Outdated or inapplicable authorities & Temporal and authority-aware retrieval. & Outdated-law error rate and applicability accuracy. \\
Incomparable research results & Shared benchmarks and reproducible protocols. & Cross-system results under common data and evaluation rules. \\
Limited auditability & Preserve provenance and traceable evidence links. & Citation resolution, provenance completeness, and audit success rate. \\
\bottomrule
\end{tabularx}
\end{table*}

These research directions are interdependent. Knowledge graphs can support authority-aware retrieval, but
their usefulness depends on accurate source extraction. Multilingual retrieval can improve access, but its
outputs still require jurisdictional filtering. Structure-aware indexing can preserve context, while
claim-level verification can assess whether the generated answer uses that context correctly. Agentic
workflows can coordinate these capabilities, but only reproducible evaluation can establish whether the
resulting system is more dependable. The most promising research agenda is therefore an integrated one:
knowledge-graph-grounded, multilingual, structure-aware legal retrieval with explicit evidence
verification and reproducible evaluation.

\subsection{Concluding synthesis}

The evolution of legal AI demonstrates a shift from systems focused primarily on rules, cases, and legal
language understanding toward architectures that retrieve external evidence and coordinate more complex
legal information workflows. Yet access to evidence is not equivalent to legal correctness. Dependable
systems must preserve the structure and provenance of legal sources, recognize jurisdictional and temporal
constraints, support multilingual queries, and verify that generated claims are adequately supported.
Three priorities emerge from this survey:
\begin{enumerate}[leftmargin=1.4em, itemsep=2pt, topsep=2pt]
\item represent legal knowledge more faithfully through hierarchical indexing, explicit relationships, and
provenance-aware knowledge graphs;
\item make evidence retrieval more accessible and reliable through hybrid and multilingual retrieval,
contextual ranking, and jurisdiction-aware filtering;
\item make generated answers more verifiable through claim-level evidence checks, authority validation,
appropriate abstention, and reproducible expert evaluation.
\end{enumerate}

For Pakistani legal AI, these priorities provide a concrete path toward systems that connect users'
questions to applicable legal evidence without treating semantic similarity or fluent generation as proof
of correctness. More broadly, they suggest that progress in legal AI should be measured not only by what a
model can generate, but by whether its answers can be traced, evaluated, and justified through the legal
sources on which they depend.
